\section{Results}

\subsection{Structured outputs and parse stability}

Both MedGemma-4B-IT and Qwen2.5-VL-7B produced structurally valid JSON for all 1,000 NIH pilot studies after the prompt and parser were finalised:

\begin{table}[t]
\centering
\small
\begin{tabular}{lll}
\hline
Model & Samples & Parse failures \\
MedGemma-4B-IT & 1000 & 0 \\
Qwen2.5-VL-7B & 1000 & 0 \\
\hline
\end{tabular}
\end{table}


However, parse stability did not imply reliable image understanding.

\subsection{Overall NIH performance and calibration}

Macro metrics across the eight findings are shown in Table 1.

\begin{table}[t]
\centering
\small
\begin{tabular}{lllllll}
\hline
Model & Split & Macro AUROC & Macro ECE & Macro Brier & Macro AURC & Mean coverage @5\% risk \\
MedGemma & internal\_train & 0.647 & 0.133 & 0.122 & 0.120 & 0.733 \\
MedGemma & source\_val & 0.636 & 0.154 & 0.124 & 0.137 & 0.713 \\
MedGemma & internal\_test & 0.612 & 0.232 & 0.198 & 0.228 & 0.357 \\
Qwen2.5-VL & internal\_train & 0.500 & 0.088 & 0.087 & 0.071 & 0.662 \\
Qwen2.5-VL & source\_val & 0.502 & 0.070 & 0.070 & 0.061 & 0.662 \\
Qwen2.5-VL & internal\_test & 0.509 & 0.153 & 0.154 & 0.117 & 0.398 \\
\hline
\end{tabular}
\end{table}


MedGemma showed moderate but non-trivial discrimination (macro AUROC 0.61--0.65) and clear calibration degradation on \texttt{internal\_test} (macro ECE 0.232). Qwen remained close to chance discrimination (macro AUROC $\approx$ 0.50), despite acceptable-looking calibration and selective metrics.

\subsection{Qwen is structurally valid but image-invariant}

On 300 NIH images, MedGemma changed many predictions under counterfactual image conditions:

\begin{table}[t]
\centering
\small
\begin{tabular}{lll}
\hline
Condition & MedGemma image sensitivity & Qwen image sensitivity \\
blank gray & 1.000 & 0.0033 \\
strong blur & 0.947 & 0.0033 \\
shuffled image & 0.877 & 0.0033 \\
\hline
\end{tabular}
\end{table}


MedGemma’s mean absolute probability change was 0.138--0.210. Qwen’s was 0.0006, effectively zero.

Qwen also predicted \texttt{no\_finding = positive} for 988/1000 studies and had macro AUROC $\approx$ 0.50 across all eight findings. Its apparently low selective risk and high selective accuracy therefore reflected a degenerate prediction strategy rather than reliable image interpretation.

\textbf{Finding:} Structured output validity, low selective risk, and high selective accuracy can coexist with image invariance.

\subsection{Shift effects are heterogeneous and non-monotonic}

MedGemma macro AUROC and ECE under shift are shown in Table 2.

\begin{table}[t]
\centering
\small
\begin{tabular}{llll}
\hline
Condition & Severity & Macro AUROC & Macro ECE \\
clean & none & 0.620 & 0.163 \\
blur & mild & 0.612 & 0.166 \\
blur & moderate & 0.603 & 0.160 \\
blur & severe & 0.615 & 0.137 \\
noise & mild & 0.675 & 0.149 \\
noise & moderate & 0.704 & 0.151 \\
noise & severe & 0.520 & 0.142 \\
JPEG & mild & 0.591 & 0.167 \\
JPEG & moderate & 0.710 & 0.149 \\
JPEG & severe & 0.621 & 0.140 \\
brightness & mild & 0.614 & 0.175 \\
brightness & moderate & 0.600 & 0.158 \\
brightness & severe & 0.597 & 0.164 \\
contrast & mild & 0.670 & 0.165 \\
contrast & moderate & 0.620 & 0.163 \\
contrast & severe & 0.617 & 0.164 \\
resolution & mild & 0.633 & 0.148 \\
resolution & moderate & 0.705 & 0.134 \\
resolution & severe & 0.581 & 0.104 \\
\hline
\end{tabular}
\end{table}


Image sensitivity generally increased with severity (e.g., noise 0.53 $\rightarrow$ 0.70 $\rightarrow$ 0.79; resolution 0.37 $\rightarrow$ 0.60 $\rightarrow$ 0.71), but AUROC did not follow a monotonic degradation pattern. Several shifts even improved calibration.

Paired macro bootstrap tests on MedGemma showed:

\begin{itemize}
\item no shift produced a significant AUROC decrease;
\item JPEG moderate: $\Delta$AUROC = +0.069, p = 0.05;
\item resolution moderate: $\Delta$AUROC = +0.068, p = 0.10;
\item significant ECE improvements for blur severe, JPEG moderate/severe, noise mild/severe, and resolution mild/moderate/severe;
\item only brightness mild showed a significant ECE increase ($\Delta$ECE = +0.010, p = 0.04).
\end{itemize}

\textbf{Finding:} Distribution shift does not produce a simple monotonic degradation of CXR VLM reliability. Shift response is condition-specific and must be evaluated jointly with image reliance and calibration.

\subsection{Calibration}

Per-finding temperature scaling fitted on \texttt{source\_val} did not change discrimination. It produced mixed calibration effects:

\begin{itemize}
\item \texttt{no\_finding} ECE improved substantially: MedGemma 0.363 $\rightarrow$ 0.093;
\item most other findings showed little improvement or slight worsening;
\item Qwen’s \texttt{no\_finding} ECE improved from 0.583 to 0.258, but AUROC stayed at chance.
\end{itemize}

These results suggest that global or per-label temperature scaling is insufficient for multi-label CXR reliability.

\subsection{Selective prediction}

On pooled \texttt{internal\_test} label-instances, MedGemma achieved:

\begin{verbatim}
global temperature: 5.67
AURC: 0.201
coverage@5% risk: 0.002
coverage@10% risk: 0.006
selective accuracy@80% coverage: 0.680
\end{verbatim}

Qwen achieved:

\begin{verbatim}
global temperature: 5.34
AURC: 0.095
coverage@5% risk: 0.000
coverage@10% risk: 0.897
selective accuracy@80% coverage: 0.901
\end{verbatim}

Qwen’s selective metrics appear strong but are invalid as clinical evidence because its predictions are image-invariant and dominated by \texttt{no\_finding}. This demonstrates that selective risk must be conditioned on image reliance.

\subsection{CheXpert-plus-RRG external report evaluation}

On 264 RRG images, MedGemma produced:

\begin{verbatim}
overall hallucination rate: 0.135
overall negation-failure rate: 0.006
overall omission rate: 0.122
overall uncertain rate: 0.090
\end{verbatim}

Finding-level results:

\begin{table}[t]
\centering
\small
\begin{tabular}{lll}
\hline
Finding & Hallucination & Omission \\
atelectasis & 0.443 & 0.057 \\
cardiomegaly & 0.449 & 0.003 \\
consolidation & 0.019 & 0.127 \\
edema & 0.051 & 0.166 \\
no\_finding & 0.083 & 0.252 \\
pleural\_effusion & 0.019 & 0.213 \\
pneumonia & 0.019 & 0.054 \\
pneumothorax & 0.000 & 0.105 \\
\hline
\end{tabular}
\end{table}


The high atelectasis/cardiomegaly hallucination rates may partly reflect model over-calling and limitations of the rule-based evidence extractor. On the first 50 RRG images, original-vs-blank status flip rate was 0.4525, indicating that MedGemma still depends on the image in this external set.

\subsection{Summary of empirical findings}

\begin{enumerate}
\item MedGemma is image-dependent, but its calibration and selective reliability are weak.
\item Qwen is structurally valid but image-invariant and near chance on findings.
\item Image-quality shift impacts prediction stability but does not monotonically degrade macro AUROC or ECE.
\item Selective metrics can look favourable for an image-invariant model.
\item RRG external report evaluation exposes both hallucination and omission behaviour.
\end{enumerate}

Together these results motivate a shift-aware selective trust framework that evaluates calibration, uncertainty, and visual evidence reliance jointly rather than relying on accuracy or selective risk alone.
